\chapter{安全基础}
\label{ch:security}
\chapterlead{安全不是一套工具，而是一组取舍：攻击面越小、权限越窄、暴露越少，出问题的概率越低。本章给出个人服务器上性价比最高的六项措施。}

\section{起步：先看暴露面}

\begin{shellsession}
$ ss -tlnp
$ sudo nft list ruleset
$ systemctl list-units --type=service --state=running
$ last -n 10
$ journalctl -u sshd | grep -c "Failed password"
\end{shellsession}

五条命令回答：开了哪些端口、防火墙规则是什么、跑了哪些服务、谁登录过、SSH 被猜了多少次。最后一项通常是个大数字，它说明公网上的扫描是持续不断的。

\section{六项措施}

\begin{booktable}{\textwidth}{J{0.815} J{1.185}}
\toprule
措施 & 挡住什么 \\
\midrule
\endhead
只开必要端口 & 扫描与自动化攻击 \\
密钥登录、禁用密码 & 暴力猜测 \\
最小权限用户 & 一旦被入侵的扩散范围 \\
自动更新与重建 & 已知漏洞的长期存在 \\
备份与异地副本 & 勒索与误删 \\
日志与告警 & 长期潜伏的异常 \\
\bottomrule
\end{booktable}

前两项在第\ref{ch:server-access}章做过，第三项在服务器部署时落实，第四项在第\ref{ch:upgrade-rollback}章。本章补上后两项的具体做法。

\section{自动更新}

\begin{nixcode}
system.autoUpgrade = {
  enable = true;
  allowReboot = false;
  dates = "04:00";
  flake = "github:weiwen/srv1";
};
\end{nixcode}

\texttt{allowReboot = false} 表示只切换配置、不自动重启。内核与部分库的更新需要重启才生效，因此还要定期手工重启，或设置为在指定时间窗口自动重启。

自动更新的风险是新版本引入回归。折中做法：自动更新到新代但不自动重启，观察一天后再重启；出问题时用旧代回滚。

\section{fail2ban 与 SSH 防护}

\begin{nixcode}
services.fail2ban = {
  enable = true;
  maxretry = 3;
  bantime = "1h";
  jails.sshd.settings = {
    enabled = true;
    port = "ssh";
  };
};
\end{nixcode}

连续三次失败的来源地址被封一小时。它挡不住分布式扫描，但能显著减少日志噪音。真正的防线仍然是密钥登录。

\begin{shellsession}
$ sudo fail2ban-client status sshd
$ journalctl -u fail2ban -n 20
\end{shellsession}

\section{密钥与秘密}

三条规则：不入库、不落盘到非加密位置、不共用。

\begin{nixcode}
# 运行时注入，不进入 store
systemd.services.api.serviceConfig.EnvironmentFile = "/run/secrets/api.env";
\end{nixcode}

\begin{shellsession}
$ ls -l /run/secrets/
$ sudo install -m 600 /dev/null /run/secrets/api.env
\end{shellsession}

\texttt{/run} 是内存文件系统，重启即清空，由部署流程重新写入。把秘密放在这里，可以避免它出现在 store、Git 历史与备份中。需要更强的方案时，可用 \texttt{sops-nix} 或 \texttt{agenix} 把加密后的秘密放进仓库，开机时用主机密钥解密。

\section{审计与回顾}

安全设置会随时间腐化：临时开放的端口忘了关，测试账户忘了删。

\begin{shellsession}
$ sudo nix-store --verify --check-contents
$ systemctl list-units --state=failed
$ grep -rn "PasswordAuthentication" /etc/nixos/
$ sudo awk -F: '$3>=1000 && $3<65534 {print $1}' /etc/passwd
\end{shellsession}

四条检查分别针对：store 完整性、失败的服务、配置里是否又打开了密码登录、有哪些普通用户。每月跑一次，十分钟内能完成。

\section{练习}

\begin{enumerate}
  \item 用五条命令列出本机的暴露面，逐项判断是否必要。
  \item 统计最近一周 SSH 失败登录次数，说明这个数字的含义。
  \item 配置 fail2ban，用错误的密码尝试三次，确认地址被封。
  \item 把一项服务秘密改为从 \texttt{/run/secrets} 读取，确认它不在 store 里。
  \item 按四项检查做一次月度审计，写出发现的问题。
\end{enumerate}

\begin{checkpoint}
本章结束时，应能：
\begin{itemize}
  \item 用五条命令评估一台服务器的暴露面；
  \item 说明六项措施各自挡住什么风险；
  \item 配置自动更新并解释为何默认不自动重启；
  \item 用 fail2ban 降低扫描噪音，知道它的局限；
  \item 按三条规则处理密钥与秘密，并做月度审计。
\end{itemize}
\end{checkpoint}
